% Generated draft: requires mathematical and source review.
\documentclass[11pt]{article}
\usepackage[margin=25mm]{geometry}
\usepackage{fontspec}
\setmainfont{Latin Modern Roman}
\usepackage{amsmath,amssymb,mathtools,unicode-math}
\setmathfont{Latin Modern Math}
\usepackage{enumitem,needspace,xurl,hyperref}
\setlength{\parindent}{0pt}
\setlength{\parskip}{5pt}
\setlength{\emergencystretch}{3em}
\newcommand{\N}{\mathbb{N}}
\newcommand{\Z}{\mathbb{Z}}
\newcommand{\Q}{\mathbb{Q}}
\newcommand{\R}{\mathbb{R}}
\newcommand{\C}{\mathbb{C}}
\newcommand{\F}{\mathbb{F}}
\newcommand{\A}{\mathbb{A}}
\newcommand{\PP}{\mathbb{P}}
\newcommand{\E}{\mathbb{E}}
\newcommand{\eps}{\varepsilon}
\newcommand{\dd}{\,\mathrm{d}}
\begin{document}
\section*{UnsolvedMath ID: 30001869}
\textbf{Source problem code:} catalogue entry
\subsection{Definitions and mathematical statement}
A Newton--Okounkov body requires both a divisor or graded linear series and a valuation, often supplied by an admissible flag. For example, under the usual complete-linear-series convention, let \(X\) have dimension \(m\), let \(D\) be a big Cartier divisor, and let
\[
Y_\bullet:\quad X=Y_0\supset Y_1\supset\cdots\supset Y_m=\{x\}
\]
be a flag of irreducible subvarieties with \(\operatorname{codim}_X(Y_r)=r\), each nonsingular at \(x\). The flag valuation \(\nu_{Y_\bullet}(s)\) records successive orders of vanishing along the flag, after dividing out the preceding vanishing orders and restricting. In this convention,
\[
\Delta_{Y_\bullet}(X,D)
=\overline{\operatorname{conv}}\left(
\bigcup_{k\geq 1}
\left\{\frac{\nu_{Y_\bullet}(s)}{k}:0\ne s\in H^0(X,\mathcal O_X(kD))\right\}
\right)\subseteq\R^m.
\]
These are explanatory conditional definitions; the supplied record does not specify that this divisor convention or these hypotheses are the intended ones.

The recorded question concerns a Bott--Samelson variety \(X=R(i_1,\ldots,i_l)\) associated with \(\operatorname{GL}_n\). It asks whether its Newton--Okounkov body for the source's ``indicated flag'' equals
\[
P=A_{i_1}T_{v_1}A_{i_2}\cdots T_{v_{l-1}}A_{i_l}(p)
\]
for suitable parameters
\[
p\in\R^d,\qquad v_j\in\R^{d_{i_j}}\quad(1\leq j\leq l-1).
\]
Here \(A_i\), \(T_v\), \(d\), and \(d_i\) are retained exactly as source notation, without assigned definitions. The supplied record omits the indicated flag, the divisor or graded linear series, and the definitions and domains of these operators. It also does not specify the ground field or the hypotheses on the index sequence. Consequently, the displayed equality records the proposed target, but the available information does not determine a fully quantified mathematical conjecture or identify the two sides in a common ambient space.
\subsection{Short English statement}
For the Bott--Samelson variety \(R(i_1,\ldots,i_l)\), the source asks whether a specified Newton--Okounkov body can be obtained by the displayed sequence of operators applied to a suitable point. The flag, divisor data, operators, and dimensions needed to make this equality precise are absent from the supplied record. This is a statement-only draft requiring human mathematical review.
\subsection{Sources}
\begin{itemize}
\item Local UnsolvedMath registry, permanent ID \(30001869\), source problem code \texttt{catalogue entry}, title: Newton--Okounkov Bodies of Bott--Samelson Varieties. The record belongs to the Oberwolfach Reports Open Problems collection and lists proposed year \(2011\); no proposer is named.
\item Supplied source citation: \emph{Mini-Workshop: New Developments in Newton-Okounkov Bodies (2012)}, Oberwolfach Reports, pp.~2327--2363, DOI \texttt{10.4171/owr/2011/41}. \url{https://doi.org/10.4171/owr/2011/41}. The background describes Oberwolfach Reports as documenting workshops and their open-problem discussions.
\item Supplied literature source, labeled ``original OWR report'': \url{https://doi.org/10.4171/owr/2011/41}.
\item Supplied literature source: \emph{Okounkov bodies and toric degenerations}, \url{https://arxiv.org/abs/1001.4566}.
\item Literature assessment dated exactly \texttt{2026-08-21}: the supplied curation labels the problem open and describes the Bott--Samelson equality as a conjecture supported by computations for length three, with special cases in the available literature but no stated general result. This is supplied dated curation, not an independently verified assessment; time-sensitive status and cited claims require independent verification.
\item The registry labels the statement \texttt{corrected\_verified} and reports that its mathematical content was checked against the official Oberwolfach report, with notation, context, and question form normalized without changing the source-level claim. This is a supplied verification claim; the cited sources have not been independently consulted for this draft.
\end{itemize}
\end{document}
